% GENERATED FILE. Edit canonical lessons, then run npm run generate:books.
% canonical-source-hash: fnv1a64:89ae1d1f0d355901
% canonical-lessons: FR-C286-sondage, FR-C286-bilan, FR-C286-hausse, FR-C286-baisse, FR-C286-statistique

\chapter{Les chiffres}
\label{ch:fr-les-chiffres}
\hlchaptermodality{286}

\begin{chapteropening}
\textbf{By the end of this chapter:} \emph{I can report figures: a poll, a final count, a rise, a fall and a statistic.}

Complete the last lesson of chapter 286: I can report figures: a poll, a final count, a rise, a fall and a statistic.
\end{chapteropening}

\section[le sondage]{le sondage — the poll, the survey}

\label{lesson:FR-C286-sondage}



\begin{quote}
Before the new one: say the French for censorship, then the French for the controversy.
\end{quote}

\subsection*{You'll want to know: le sondage}
\textbf{le sondage} — \textquotedblleft{}the poll, the survey\textquotedblright{}.

\begin{quote}
Un sondage montre que les jeunes lisent moins.
\end{quote}

\begin{quote}
\emph{A poll shows that young people read less.}
\end{quote}

\emph{Un sondage d'opinion} — an opinion poll.

\begin{cousinweb}
\emph{sondage} is built on \emph{sonder} \textquotedblleft{}to sound, to probe\textquotedblright{}, from \emph{sonde} \textquotedblleft{}a sounding line\textquotedblright{}, which is usually traced to Old English \emph{sund} \textquotedblleft{}water, sea\textquotedblright{}; the origin is not fully settled (EtymDB lists no parent).
\end{cousinweb}

\begin{grammarlens}[title={What you've built}]
One more word for telling, explaining and describing at B1.
\end{grammarlens}

\subsection*{Your turn}
\begin{itemize}
\raggedright
  \item \emph{Say it:} \emph{le sondage}
  \item \emph{Say it:} \emph{le sondage}, once more
  \item \emph{Say it:} a sentence of your own with \emph{le sondage}
  \item \emph{Recall:} say the French for censorship, then the French for the controversy, then say \emph{le sondage} again
\end{itemize}

\begin{tcolorbox}[breakable,colback=teal!4,colframe=teal!35!black,title={Before you move on}]
What does le sondage mean? (\textquotedblleft{}The poll, the survey\textquotedblright{}.) Say it once more.
\end{tcolorbox}
\section[le bilan]{le bilan — the final count, the outcome; the review}

\label{lesson:FR-C286-bilan}



\begin{quote}
Before the new one: say the French for the controversy, then the French for the poll.
\end{quote}

\subsection*{You'll want to know: le bilan}
\textbf{le bilan} — \textquotedblleft{}the final count, the outcome; the review\textquotedblright{}.

\begin{quote}
Le bilan de l'année est très positif.
\end{quote}

\begin{quote}
\emph{The outcome of the year is very positive.}
\end{quote}

After a disaster, \emph{le bilan} is the number of dead and injured. \emph{Faire le bilan} — to take stock.

\begin{cousinweb}
\emph{bilan} was borrowed from Italian \emph{bilancio} \textquotedblleft{}balance\textquotedblright{} (EtymDB), built on Late Latin \emph{bilanx} \textquotedblleft{}with two scales\textquotedblright{} (Lewis \& Short).
\end{cousinweb}

\begin{grammarlens}[title={What you've built}]
One more word for telling, explaining and describing at B1.
\end{grammarlens}

\subsection*{Your turn}
\begin{itemize}
\raggedright
  \item \emph{Say it:} \emph{le bilan}
  \item \emph{Say it:} \emph{le bilan}, once more
  \item \emph{Say it:} a sentence of your own with \emph{le bilan}
  \item \emph{Recall:} say the French for the controversy, then the French for the poll, then say \emph{le bilan} again
\end{itemize}

\begin{tcolorbox}[breakable,colback=teal!4,colframe=teal!35!black,title={Before you move on}]
What does le bilan mean? (\textquotedblleft{}The final count, the outcome; the review\textquotedblright{}.) Say it once more.
\end{tcolorbox}
\section[la hausse]{la hausse — the rise, the increase}

\label{lesson:FR-C286-hausse}



\begin{quote}
Before the new one: say the French for the poll, then the French for the final count.
\end{quote}

\subsection*{You'll want to know: la hausse}
\textbf{la hausse} — \textquotedblleft{}the rise, the increase\textquotedblright{}.

\begin{quote}
Les prix sont en hausse.
\end{quote}

\begin{quote}
\emph{Prices are going up.}
\end{quote}

\emph{En hausse} — going up.

\begin{cousinweb}
\emph{hausse} is built on \emph{hausser} \textquotedblleft{}to raise\textquotedblright{} (EtymDB), which is inherited from Vulgar Latin \emph{altiāre}, built on Latin \emph{altus} \textquotedblleft{}high\textquotedblright{}.
\end{cousinweb}

\begin{grammarlens}[title={What you've built}]
One more word for telling, explaining and describing at B1.
\end{grammarlens}

\subsection*{Your turn}
\begin{itemize}
\raggedright
  \item \emph{Say it:} \emph{la hausse}
  \item \emph{Say it:} \emph{la hausse}, once more
  \item \emph{Say it:} a sentence of your own with \emph{la hausse}
  \item \emph{Recall:} say the French for the poll, then the French for the final count, then say \emph{la hausse} again
\end{itemize}

\begin{tcolorbox}[breakable,colback=teal!4,colframe=teal!35!black,title={Before you move on}]
What does la hausse mean? (\textquotedblleft{}The rise, the increase\textquotedblright{}.) Say it once more.
\end{tcolorbox}
\section[la baisse]{la baisse — the drop, the fall}

\label{lesson:FR-C286-baisse}



\begin{quote}
Before the new one: say the French for the final count, then the French for the rise.
\end{quote}

\subsection*{You'll want to know: la baisse}
\textbf{la baisse} — \textquotedblleft{}the drop, the fall\textquotedblright{}.

\begin{quote}
Le nombre de touristes est en baisse.
\end{quote}

\begin{quote}
\emph{The number of tourists is falling.}
\end{quote}

\emph{En baisse} — going down; the verb is \emph{baisser}.

\begin{cousinweb}
\emph{baisse} is built on \emph{baisser} \textquotedblleft{}to lower\textquotedblright{}, inherited from Vulgar Latin \emph{bassiāre}, built on Late Latin \emph{bassus} \textquotedblleft{}low\textquotedblright{} (EtymDB).
\end{cousinweb}

\begin{grammarlens}[title={What you've built}]
One more word for telling, explaining and describing at B1.
\end{grammarlens}

\subsection*{Your turn}
\begin{itemize}
\raggedright
  \item \emph{Say it:} \emph{la baisse}
  \item \emph{Say it:} \emph{la baisse}, once more
  \item \emph{Say it:} a sentence of your own with \emph{la baisse}
  \item \emph{Recall:} say the French for the final count, then the French for the rise, then say \emph{la baisse} again
\end{itemize}

\begin{tcolorbox}[breakable,colback=teal!4,colframe=teal!35!black,title={Before you move on}]
What does la baisse mean? (\textquotedblleft{}The drop, the fall\textquotedblright{}.) Say it once more.
\end{tcolorbox}
\section[la statistique]{la statistique — the statistic; statistics}

\label{lesson:FR-C286-statistique}



\begin{quote}
Before the new one: say the French for the rise, then the French for the drop.
\end{quote}

\subsection*{You'll want to know: la statistique}
\textbf{la statistique} — \textquotedblleft{}the statistic; statistics\textquotedblright{}.

\begin{quote}
Les statistiques montrent une baisse des accidents.
\end{quote}

\begin{quote}
\emph{The statistics show a fall in accidents.}
\end{quote}

Plural in the sense of figures; \emph{la statistique} is also the field of study.

\begin{cousinweb}
\emph{statistique} was borrowed from German \emph{Statistik} (EtymDB), built on Latin \emph{status} \textquotedblleft{}state\textquotedblright{}, from \emph{stāre} \textquotedblleft{}to stand\textquotedblright{}: first the science of the state.
\end{cousinweb}

\begin{grammarlens}[title={What you've built}]
One more word for telling, explaining and describing at B1.
\end{grammarlens}

\subsection*{Your turn}
\begin{itemize}
\raggedright
  \item \emph{Say it:} \emph{la statistique}
  \item \emph{Say it:} \emph{la statistique}, once more
  \item \emph{Say it:} a sentence of your own with \emph{la statistique}
  \item \emph{Recall:} say the French for the rise, then the French for the drop, then say \emph{la statistique} again
\end{itemize}

\begin{tcolorbox}[breakable,colback=teal!4,colframe=teal!35!black,title={Before you move on}]
What does la statistique mean? (\textquotedblleft{}The statistic; statistics\textquotedblright{}.) Say it once more.
\end{tcolorbox}
